\documentclass{article}

\usepackage[utf8]{inputenc}
\usepackage{authblk}
\usepackage{hyperref}
\usepackage{enumitem}

\usepackage{array}


%\setlength{\oddsidemargin}{-18mm}\setlength{\evensidemargin}{-18mm}


%\setlength{\oddsidemargin}{-5mm} %\setlength{\evensidemargin}{-5mm}   %%
\setlength{\topmargin}{-20mm}
%\setlength{\textwidth}{194mm} \setlength{\textheight}{270mm}


\usepackage{tikz}
\usepackage{tikz-cd}
\usetikzlibrary{arrows}
\tikzcdset{arrow style=tikz,
diagrams={>={Stealth[round,length=4pt,width=6pt,inset=3pt]}}}

\usepackage{amsthm}

\usepackage{amsmath} 
\usepackage{amssymb}
\usepackage{bbold}
\usepackage{amsfonts}
\usepackage{faktor}
\usepackage[T1]{fontenc}
\usepackage{faktor}
\usepackage{xfrac}
\usepackage{dsfont}
\usepackage{color}
 \usepackage{todonotes}
\usepackage[english]{babel}
\usepackage{mathtools}

 
\newtheorem{theorem}{Theorem}
\newtheorem{definition}[theorem]{Definition}
\newtheorem{conjecture}[theorem]{Conjecture}
\newtheorem{lemma}[theorem]{Lemma}
\newtheorem{notation}[theorem]{Notation}
\newtheorem{observation}[theorem]{Observation}
\newtheorem{proposition}[theorem]{Proposition}
\newtheorem{property}[theorem]{Property}
\newtheorem{example}[theorem]{Example}
\newtheorem{remark}[theorem]{Remark}
\newtheorem{corollary}[theorem]{Corollary}

\newtheorem*{theorem*}{Theorem}

\DeclareMathOperator{\likelihood}{likelihood}


\title{physics}
\author{ }
\date{March 2026}

\begin{document}

\maketitle

\section{Introduction}

\subsection{Beyond the Standard Model Searches}

The Standard Model (SM) of the elementary particles of modern physics and their interactions is an excellent description of the behavior of nature at the realm of high enery physics. However, it leaves some open questions to be discovered. To name a few - the nature of dark matter, the origin of neutrino masses and the matter-antimatter asymmetry of the universe. The answer to these questions is lying in the formulation of some Beyond Standard Model (BSM) physics.

In the search for BSM physics, powerful particle colliders, such as the Large Hadron Collider (LHC) at CERN, have been in use, due to their ability to poke into the high energy regime where new physics may manifest itself. While current searches based on the LHC data have failed to show any evidence for BSM physics, the generated data has been accumulated for yeara and for future analysis.

A classic analysis of a model that extends the SM consists of several human-intensive steps. Those include speculating a model, formulating predictions to where in the observable parameter space should this model differ from the SM, defining a distinctive experiment, refining cuts, simulating and more. This process may take a group of physicists several years to complete, while testing a specific hypothesis.

While the data generated by the LHC is rich and complex, most theoretically motivated BSM models have been ruled out in recent years' research. Following is a suggestion to a model-agnostic method to seach for deviations from the SM predictions, in order to direct future research efforts.

\subsection{searches for deviartions from the Standard Model}
TODO{write section}

\subsection{symmetrized}
TODO{write section}

\subsubsection{Search for Lepton Flavor Violating Processes}

One particular section of the SM particles that draws attention of searches for new physics is the lepton sector. The SM predicts leptons of three generations, or flavors - the electron, the muon and the tau, each flavor contains a lepton and a neutrino. The foundations of the SM predict three identical generations that is broken only by Yukawa interactions. We would say that there is a Lepton Flavor Universality (LFU) in the SM. However, hints such as neutrino oscillations suggest that there is physics we don't yet understand in that sector.

\subsection{In this Paper}

Our goal is to use previoulsy suggested method, the Data-Directed Paradigm (DDP), to look for hints of lepton flavor violating (LFV) evidence in simulated LHC data. To this end, we'll extend the existing method in several ways:
\begin{enumerate}
    \item We will look at k-dimensional data, in terms of physical observables of the events.
    \item We will model our knowledge of the effect of a detector and model our error in that using nuisance parameters. We'll measure its effect on the discovery probability and significance estimation.
    \item We'll proceed to run on simulated CMS datasets, rather that purely hypothetical samples
\end{enumerate}



\section{Adding Nuisance Parameters to Account for Imperfect knowledge}

We would now formulate our confrontation with the finite accuracy with which we know to model the detector. In theory, we can formulate in the same manner how our theory depends on imperfections of all sorts, but modeling this method in further parameters is left for future work. This being said, the method with which we account for detector efficiency imperfection is easily translatable to any other parameter in the model.

We would now elaborate on the structure of the likelihood. Having added nuisance parameters, with the expectance for them to be small, add another layer to how likely we think a certain hypothesis is. This is, we would consider a hypothesis in which the nuisance parameters are small more likely than one in which they are large. Of course, there is a tradeoff here to fitting the data well that we should address.

From \ref{eq:pooled_log_likelihood} and discussions about likelihood in ref (\cite{Italian_no_2_Learning_New_Physics_from_an_Imperfect_Machine}), we claim %eq 3 there

$$
\begin{aligned}
    \likelihood(\dataset{A}, \nu^j_m) = \likelihood(\dataset{A}| \nu^j_m) \cdot \likelihood(\nu^j_m)
\end{aligned}
$$

To estimate likelihood of a value of a nuisance parameter, we should assume something about the way it behaves. For simplicity and for stating the foundations of our method, let us assume a normal disribution with some $\sigma_0$ that is centered around 0, w.l.o.g. 

The likelihood of a gaussian distributed parameter is well known, as expressed By

$$
\begin{aligned} \label{eq:gaussian_likelihood}
    \likelihood(\nu^j_m) = \frac{1}{\sqrt{2 \pi \sigma_0^2}} e^{-\frac{(\nu^j_m)^2}{2 \sigma_0^2}}
\end{aligned}
$$

The logarithm of which is

$$
\begin{aligned} \label{eq:gaussian_log_likelihood}
    \llikelihood(\nu^j_m) = -\frac{1}{2} \log(2 \pi \sigma_0^2) - \frac{(\nu^j_m)^2}{2 \sigma_0^2} = - \frac{1}{2 \sigma_0^2} (\nu^j_m)^2 + C
\end{aligned}
$$

Leaving constants out, we incorporate this into the already existing "nuisance-clean" estimation of the likelihood of a dataset. Using the incorporation of nuisance parameters in our equation \ref{eq:efficiency_j_explicit_form}, \ref{eq:def_weights} and plugging it and \ref{eq:gaussian_log_likelihood} into \ref{eq:symmetrized_test_statistics} to get

$$
\begin{aligned} \label{eq:def_test_statistic_with_nuisance}
t_{\dataset{A}+\dataset{B}}(\dataset{A}) &= -2 \log \big( \likelihood(\dataset{A}| \nu^j_m) \cdot \likelihood(\nu^j_m) \big) = -2 \log \likelihood(\dataset{A}|\nu^j_m) + \sum_{i=0,j=0}^{N_\reco{A}, k} \log \likelihood(\nu^j_m) \\
 &= -2 \cdot \min_{f(x)}
\left[
    - \frac{N_{\reco{A}}}{N_{\reco{A}} + N_{\reco{B}}}
    \sum_{x_i \in \reco{A} \cup \reco{B}} \frac{1}{\theta(x_i)} (e^{f(x_i)} - 1)
    + \sum_{x_i \in \reco{A}} \frac{1}{\theta(x_i)} f(x_i)
    - \sum_{\nu^j} \frac{1}{2 \sigma_0^2} (\nu^j_m)^2
\right]
\end{aligned}
$$

and equivalently for $t_{\dataset{A}+\dataset{B}}(\dataset{B})$.

This being said, we modify the loss function accordingly to

$$
\begin{aligned} \label{eq:loss_function_with_nuisance}
L[f, \nu^j] &=
    - \frac{N_{\reco{A}}}{N_{\reco{A}} + N_{\reco{B}}} \frac{1}{\theta(x_i)}
    \big[
        \sum_{x_i \in \reco{A} \cup \reco{B}}
        (1-y_i)(e^{f(x_i)} - 1)
        - y_i f(x_i)
    \big]
    - \sum_{\nu^j} \frac{1}{2 \sigma_0^2} (\nu^j)^2
\end{aligned}
$$

Note that both the weights (through $\theta$'s) and the nuisance likelihood terms depend on the nuisance parameters, in ways that contribute opposingly. Maximizing this value using our NN gradient descent method would find the most probable solution for these constriants.


\section{Error Propagation in Detector Efficiency}
In practical scenarios, the detector efficiency \( \eta(x) \) is estimated and carries uncertainties \( \delta \eta(x) \). These propagate into the weights, reconstructed event counts, the test statistic, and the training loss:

\subsection{Uncertainty in Event Weights}
The weights depend inversely on efficiency:
\[
w_i = \frac{1}{\eta(x_i)} \implies \delta w_i \approx \frac{\delta \eta(x_i)}{\eta(x_i)^2}
\]
where we have taken the first-order derivative \(\partial w / \partial \eta = -1/\eta^2\) and noted that the sign is irrelevant for variances.

\subsection*{Proposition 2 (Variance of reconstructed counts).}
If efficiency uncertainties for different events are uncorrelated,
\[
(\delta \dot{N}_{\mathbf{A}})^2 = \sum_{x_i \in \tilde{\mathbf{A}}} \left( \frac{\delta \eta(x_i)}{\eta(x_i)^2} \right)^2,
\]
and similarly for \( \mathbf{B} \).  
\emph{Proof.} Using the standard error propagation formula for independent variables \(\eta(x_i)\):
\[
(\delta \dot{N}_{\mathbf{A}})^2 = \sum_i \left( \frac{\partial \dot{N}_{\mathbf{A}}}{\partial \eta(x_i)} \delta \eta(x_i) \right)^2 = \sum_i \left( -\frac{1}{\eta(x_i)^2} \delta \eta(x_i) \right)^2.
\]
\(\square\)

\subsection{Uncertainty Propagation into the Test Statistic}
The test statistic depends on \( \dot{N}_{\mathbf{A}} \), \( \dot{N}_{\mathbf{B}} \), and functions \( f(x), g(x) \) trained on weighted data. To first order, the propagated uncertainty is
\[
(\delta t)^2 \approx \left( \frac{\partial t}{\partial \dot{N}_{\mathbf{A}}} \delta \dot{N}_{\mathbf{A}} \right)^2 + \left( \frac{\partial t}{\partial \dot{N}_{\mathbf{B}}} \delta \dot{N}_{\mathbf{B}} \right)^2 + \cdots,
\]
where higher-order contributions and training uncertainties can be estimated numerically.

\subsection{Incorporating Uncertainty into the Loss Function}
Uncertainty in weights also propagates into the loss:
\[
\delta L = \sqrt{
\sum_{i} \left( \frac{\partial L}{\partial w_i} \delta w_i \right)^2
} = \sqrt{
\sum_{i} \left( \frac{\partial L}{\partial w_i} \frac{\delta \eta(x_i)}{\eta(x_i)^2} \right)^2
}.
\]
A practical approach is to add a regularization term:
\[
L_{\mathrm{total}} = L + \lambda \cdot \delta L,
\]
where \(\lambda\) controls the tradeoff between fitting the data and accounting for efficiency uncertainty.

\subsection{Proposition 3 (Bootstrap consistency).}
Let \(t\) be any statistic computed from the efficiency-weighted data. If \(\hat{F}\) is the empirical distribution of \((x_i, w_i)\), then the bootstrap variance estimator
\[
\hat{\sigma}^2_{\mathrm{boot}}(t) = \frac{1}{B-1} \sum_{j=1}^B \left( t^{(j)} - \bar{t} \right)^2
\]
converges in probability to \(\mathrm{Var}(t)\) as \(B \to \infty\), provided \(t\) is a smooth functional of \(F\). This follows from standard nonparametric bootstrap theory. \(\square\)

\section{Formalism}

\subsection{Building an Observed Dataset}

During a HEP experiment, large amount of collision events occur inside the detector. The original number of such events in a dataset is denoted by $\true{N}$.

When detected, some $j$ physical attributes of each event are measured, e.g., energy, momentum components and angles. Let each event \( \bar{x_i} \in \mathbb{R}^k \equiv \{\phi^j_i\}\) be a point in \( k \)-dimensional feature space. Each dataset \(\mathbf{A} \in \mathbb{R}^{\true{N} \times k}\) consists of such events.

Due to limited detector efficiency and selection effects, not all events that occur inside the detector are observed. Let $\eta(\bar{x})$ be the probability with which an event is detected, and hence appears in the observed dataset.

We define the filtering process \(\detector(\eta, \{\bar{x_i}\})\) that stochastically removes events:
\[
\detector(\eta, \{\bar{x_i}\}) \equiv \{ \bar{x_i} \mid \text{each } \bar{x_i} \text{ is retained with probability } \eta(\bar{x_i}) \}.
\]

The observed (filtered) dataset is then:
\[
\meas{A} = \detector(\eta, \dataset{A})
\]

with $|\meas{A}| = N_\meas{A}$, the observed number of events.

Our first goal is to estimate \( \true{N} \) from the observed \( N_\meas{A} \), correcting for detector efficiency effects and propagating uncertainties associated with the efficiency estimation.

\subsection{Modelling our Knowledge of the Detector Efficiency}

Now expressing the (true) efficiency functions as:

\begin{align*}
    \eta(\bar{x_i}) &: \mathbb{R}^k \to [0,1] \equiv \prod_{j=1}^k \eta^j(\phi_i^j)
\end{align*}

with

\begin{align} \label{eq:efficiency_j_explicit_form}
    \theta^j(\phi^j) &: \mathbb{R} \to [0,1] \equiv e^{\nu^j_m}\eta^j(\phi^j)
\end{align}

with $\theta^i$ representing our knowledge of the efficiency of each bin. This means, we model for an efficiency function that varies from our understanding of it through usage of nuisance parameters $\nu^j_m$. While considered small for our formulation, for the time being - are unconstrained.

Datasets statistics are being binned before they can be analyzed. Given that the detectable range is divided into $m^j$ bins along each dimension, we consider the efficiency of the detector to be a multiplicative constant in each. We further claim and model $\eta$ and $\theta$ to be evaluated at a per-bin basis. This gives taht we need exactly $j \times m$ nuisance parameters, $\nu^j_m$, to fully capture the uncertainty effect.

\subsection{Efficiency Compensation via Event Weights}

To reconstruct the statistics of the original datasets, we assign weights:

\begin{equation}
\label{eq:def_weights}
\text{for } x_i \in \tilde{\mathbf{A}}, \quad w_i := \frac{1}{\theta(x_i)}
\end{equation}

This gives rise to the efficiency-compensated reconstructed datasets:

\[
\dot{\mathbf{A}} := \{ (x_i, w_i) \mid x_i \in \tilde{\mathbf{A}} \}
\]

The estimated total event counts become:
\[
N_\reco{A} = \sum_{x_i \in \meas{A}} w_i \approx \true{N}
\]

Note that this expression includes our knowledge of the efficiency and our error, throught the $\eta$'s in the expression.

\subsection{Proposition (Unbiasedness of the weighted estimator).} \label{prop:unbiasedness_of_weights}
Let \(\{x_i\}_{i=1}^{\true{N}}\) be the set of true collected events, each detected independently with probability \(\eta(x_i)\). Then:
\[
\mathbb{E}\!\left[ N_\reco{A} \right] = \true{N}
\]
\emph{Proof.}
For a single event \(x_i\), the contribution to the sum is \(\frac{1}{\theta(x_i)}\) if detected and \(0\) otherwise. Its expectation is:
\begin{align*}
\eta(x_i) \cdot \frac{1}{\theta(x_i)} + (1-\eta(x_i))\cdot 0 &= \\
= e^{\nu(x_i)}\theta(x_i) \cdot \frac{1}{\theta(x_i)} &= e^{\nu(x_i)}
\end{align*}

% If we leave this formulation take the proof from https://statproofbook.github.io/P/norm-mgf.html
% It says that E(x) = e^(\mu x + 1/2 sigma^2 x^2)

Given that the expectation value of the exponentiated nuisances is 0, summing over all \(\true{N}\) events yields the result. \(\square\)


\section{Hypotheses in the Weighted Setting}
We now define the null and alternative hypotheses using reweighted densities:
\[
\begin{aligned}
\mathcal{H}_0 &: \quad
\begin{cases}
    {\dist{A}} = \dfrac{N_\reco{A}}{Z} e^{h(x)} \refdist \\
    {\dist{B}} = \dfrac{N_\reco{B}}{Z} e^{h(x) + r} \refdist
\end{cases} \\
\mathcal{H}_1 &: \quad
\begin{cases}
    {\dist{A}} = \dfrac{N_\reco{A}}{Z} e^{f(x)} \refdist \\
    {\dist{B}} = \dfrac{N_\reco{B}}{Z} e^{g(x)} \refdist
\end{cases}
\end{aligned}
\]
Here \(n_{\dataset{A}}(x)\) and \(n_{\dataset{B}}(x)\) are \emph{differential event densities} (counts per unit \(k\)-dimensional volume) rather than raw counts.


\section{Symmetrized Test Statistic with Efficiency Weights}
Following the logic in \cite{LastDDPpaper}, the symmetrized test statistic becomes:
\[
\begin{aligned} \label{eq:symmetrized_test_statistics}
t_{\dataset{A}+\dataset{B}}(\dataset{A}) &= -2 \cdot \min_{f(x)}
\left[
    - \frac{N_{\reco{A}}}{N_{\reco{A}} + N_{\reco{B}}}
    \sum_{x_i \in \reco{A} \cup \reco{B}} w_i (e^{f(x_i)} - 1)
    + \sum_{x_i \in \reco{A}} w_i f(x_i)
\right] \\
 t_{\dataset{A}+\dataset{B}}(\dataset{B}) &= -2 \cdot \min_{g(x)}
\left[
    - \frac{N_{\reco{B}}}{N_{\reco{A}} + N_{\reco{B}}}
    \sum_{x_i \in \reco{A} \cup \reco{B}} w_i (e^{g(x_i)} - 1)
    + \sum_{x_i \in \reco{B}} w_i g(x_i)
\right].
\end{aligned}
\]

\subsection*{Full rigorous proof of the above expression}
\label{sec:full-proof-t}
We present a self-contained derivation that starts from the extended (inhomogeneous) Poisson likelihoods and proceeds to the displayed convex functional. The derivation explicitly tracks constants and shows how integrals over the true (unfiltered) process become weighted sums over detected events.

\paragraph{1. Model and parametrization.}
Model the "true" events for samples $\dataset{A}$ and $\dataset{B}$ as independent inhomogeneous Poisson processes with intensities
\[
\lambda_{\dataset{A}}(x)=\true{N^{\dataset{A}}} p_{\dataset{A}}(x),\qquad
\lambda_{\dataset{B}}(x)=\true{N^{\dataset{B}}} p_{\dataset{B}}(x),
\]
where \(p_{\dataset{A}}\) and \(p_{\dataset{B}}\) are probability densities. Introduce a common reference density \(n_{\mathcal{R}}(x)\) and parametrize
\[
p_{\dataset{A}}(x)=\frac{1}{Z_\dataset{A}} e^{f(x)} \refdist,\qquad
p_{\dataset{B}}(x)=\frac{1}{Z_\dataset{B}} e^{g(x)} \refdist,
\]
with
\(Z_\dataset{A}=\int e^{f(x)} \refdist dx\) and \(Z_\dataset{B}=\int e^{g(x)} \refdist dx\) ensuring normalization. (We keep \(Z_\dataset{A},Z_\dataset{B}\) explicit; later they cancel or become irrelevant for the test statistic.)

\paragraph{2. Extended log-likelihood for each sample.}
For an inhomogeneous Poisson process with intensity \(\lambda(x)\), the log-likelihood for observed points \(\{x_i\}\) is
\[
\llikelihood(\lambda) = -\int \lambda(x)\,dx + \sum_i \log\lambda(x_i) + C
\]
where \(C\) collects terms not depending on \(\lambda\). For sample $\meas{A}$ (observed events) this gives
\[
\llikelihood_{\dataset{A}}(f) = -{N}_{\dataset{A}} + \sum_{x_i \in \dataset{A}} \log\big({N}_{\dataset{A}}\tfrac{1}{Z_\dataset{A}} e^{f(x_i)} \refdist[x_i]\big) + C_\dataset{A}
\]
Collect the \(f\)-dependent pieces:
\[
\llikelihood_{\dataset{A}}(f) = -N_{\dataset{A}} + \sum_{x_i \in \dataset{A}} f(x_i) + N_{\dataset{A}}\log{N_{\dataset{A}}} - N_{\dataset{A}}\log Z_\dataset{A} + \sum_{x_i \in \dataset{A}} \log \refdist[x_i] + C_A
\]
Analogous expression holds for sample $\dataset{B}$.

\paragraph{3. Passing to detected (filtered) data using weights.}
We do not observe the unfiltered lists, but instead the detected subsets \(\meas{A}\) and \(\meas{B}\). By the construction of weights as in (\ref{eq:def_weights}), we convert our formulation to use the reconstructed datasets \(\reco{A}\) and \(\reco{B}\). Integrals and sums over the unfiltered process are unbiasedly estimated by weighted sums over detected events (Proposition \ref{prop:unbiasedness_of_weights}). In particular,
\[
\sum_{x_i \in \dataset{A}} f(x_i) \quad\longrightarrow\quad \sum_{x_i \in \reco{A}} w_i f(x_i)
\]

while the integral term \(-\int \lambda_{\dataset{A}} = -N_{\reco{A}}\) remains explicit. Thus, working with reconstructed data, the log-likelihood for sample $\reco{A}$ (keeping only terms that may depend on \(f\)) is represented by
\[
\llikelihood_{\reco{A}}(f) \propto -N_{\reco{A}} + \sum_{x_i \in \reco{A}} w_i f(x_i) - N_\reco{A}\log Z_\reco{A} + N_\reco{A}\log N_\reco{A} + \sum_{x_i \in \reco{A}} w_i\log \refdist[x_i]
\]
The analogous expression holds for $\reco{B}$ with \(g\) and \(Z_\reco{B}\).

\paragraph{4. Pooled (symmetrized) likelihood.}
The symmetrized construction treats the pooled set of detected events \(\meas{P}=\meas{A} \cup \meas{B}\) as being drawn from a mixture of the two models. Define

\[
N_\reco{P} = N_\reco{A} + N_\reco{B},\qquad \alpha=\frac{N_\reco{A}}{N_\reco{P}}
\]

For the pooled model one can write the pooled intensity (per unit \(x\)) as

\[
\lambda_{P}(x) = N_\reco{P}\,\Big[\alpha \frac{1}{Z_\reco{A}}e^{f(x)} + (1-\alpha)\frac{1}{Z_\reco{B}}e^{g(x)}\Big] \refdist
\]

The pooled log-likelihood (for observed pooled points \(x_i\in\reco{P}\)) is

\[ \label{eq:pooled_log_likelihood}
\log\mathcal{L}_{\reco{P}}(f,g) = -N_\reco{P} + \sum_{x_i\in\reco{P}} w_i \log\Big(N_\reco{P}\big[\alpha \tfrac{1}{Z_\reco{A}}e^{f(x_i)} + (1-\alpha)\tfrac{1}{Z_\reco{B}}e^{g(x_i)}\big] \refdist[x_i] \Big) + C_{\reco{P}}
\]

Drop terms that do not depend on \(f,g\) (constants and \(\log \refdist\) terms).

\paragraph{5. Constructing the class-specific test functional.}
To obtain the test statistic that assesses whether sample A differs from the pooled reference in the manner parametrized by \(f\), set the alternative for B to be the pooled baseline. For simplicity and without loss of generality we now take the parametrization where the B-component is the baseline with \(g(x)\equiv 0\) (equivalently measure \(f\) relative to the B baseline). With this choice

\[
\log\mathcal{L}_{\reco{P}}(f) \propto -N_\reco{P} + \sum_{x_i\in\reco{P}} w_i \log\Big(\alpha \tfrac{1}{Z_\reco{A}}e^{f(x_i)} + (1-\alpha)\tfrac{1}{Z_\reco{B}}\Big).
\]

Expand the logarithm by factoring \(\tfrac{1}{Z_\reco{B}}\) out of the argument:

\[
... = \log\Big(\tfrac{1}{Z_\reco{B}}\Big) + \log\Big((1-\alpha) + \alpha\tfrac{Z_\reco{B}}{Z_\reco{A}}e^{f}\Big).
\]
The constant \(\log(1/Z_\reco{B})\) does not depend on \(f\), so discard it. Define the ratio \(\rho:=\tfrac{Z_\reco{B}}{Z_\reco{A}}\) and rewrite the remaining term as
\[
\log\Big((1-\alpha) + \alpha\rho e^{f}\Big).
\]
Now form the log-likelihood ratio (or equivalently the difference between the pooled log-likelihood and the product of the separate log-likelihoods under the null) and collect terms that depend on \(f\). The algebraic manipulations (expand and regroup the pooled-sum and the individual sums) produce two types of contributions:
\begin{enumerate}
    \item A term proportional to \(\int \alpha (\rho e^{f}-1) \refdist dx\), which, after replacing integrals with weighted sums, yields \(-\alpha\sum_{\reco{P}} w_i (e^{f(x_i)}-1)\) up to multiplicative constants and absorbed \(\rho\) factors. The appearance of \(e^{f}-1\) is a consequence of subtracting the baseline intensity.
    \item A sum over events in sample A giving \(\sum_{x_i \in \reco{A}} w_i f(x_i)\) originating from the \(\sum\log e^{f}=\sum f\) contribution of labeled-A events when comparing labeled likelihood terms.
\end{enumerate}
These two pieces are the ones that appear in the convex functional minimized to obtain the profile likelihood.

\paragraph{6. Profiling and simplification to the displayed functional.}
Carrying out the algebra carefully (expand the pooled log, subtract the baseline contributions, cancel constants and terms involving \(\log Z_A,\log Z_B\) that do not depend on \(f\) after profiling), and using the unbiased replacement of integrals by weighted sums, one obtains the functional
\[
\mathcal{F}(f) = -\alpha \sum_{x_i\in\reco{P}} w_i\big(e^{f(x_i)}-1\big) + \sum_{x_i\in\reco{A}} w_i f(x_i) + \text{(constants independent of $f$)}.
\]
The test statistic is the negative twice the maximized log-likelihood ratio (or, equivalently, \(-2\) times the minimum of \(\mathcal{F}(f)\) over allowable functions \(f\)), which yields the stated expression:
\[
t_{\dataset{A}+\dataset{B}}(\dataset{A}) = -2 \min_f \Big[ -\alpha \sum_{x_i\in\reco{P}} w_i (e^{f(x_i)}-1) + \sum_{x_i\in\reco{A}} w_i f(x_i) \Big],
\]
with \(\alpha=\dfrac{N_\reco{A}}{N_\reco{A}+N_\reco{B}}\).

\paragraph{7. Remarks on exactness and convexity.}
The steps above are algebraic and exact up to the choice of baseline parametrization (we set \(g\equiv 0\) for the baseline without loss of generality) and the unbiased substitution of integrals by weighted sums (which is exact in expectation and valid as an estimating equation). The resulting functional is convex in the pointwise values \(f(x_i)\) because the map \(f\mapsto e^{f}-1\) is convex; this justifies numerical minimization via deterministic optimization or training a neural network with the corresponding loss.

\section{Loss Function for Neural Network Training}
To learn the optimal \( f(x) \), we train a classifier using the following efficiency-weighted loss:
\[
L[f] = \sum_{x_i \in \reco{A} \cup \reco{B}} 
\left[
    (1 - y_i) \cdot \frac{N_\reco{A} w_i}{N_\reco{A} + N_\reco{B}} (e^{f(x_i)} - 1)
    - y_i w_i f(x_i)
\right]
\]
where \( y_i = 1 \) if \( x_i \in \reco{A} \), and \( y_i = 0 \) otherwise.

\section{Reproducing Previous Results}

The following results are generated using the new, improved codebase, with a subset of parameters that mimic the conditions for the previous paper \cite{LastDDPPaper}.


\end{document}
